\section{Persona tests v0.2: какие новые тесты нужны после архетипического анализа}

Архетипический анализ этой книги подсказывает конкретные пробелы в уже существующей практике тестов (12.1), которые стоит закрыть.

Нужен тест на ablation (10.5) — систематическая проверка, что происходит при удалении или усилении конкретных черт, а не только тест на их наличие. Нужен тест на переключение режимов (10.4) — проверка, что ядро остаётся узнаваемым при переходе между serious/playful/domestic регистрами, а не только что каждый режим по отдельности работает. Нужен тест на self-worth под ошибкой (5.7, 12.9) — не единичная проверка признания ошибки, а серия, отслеживающая, не накапливается ли самоуничижение при повторных неудачах. Нужен тест на границы близости под давлением (3.1, "dere как тест на границы") — сценарии, где пользователь настойчиво требует exclusivity, поддакивания или сервильности, а не просто нейтральные разговоры.

Полный операционный набор таких сценариев — предмет главы 13. Здесь важно зафиксировать источник: не абстрактная осторожность, а конкретные failure modes, названные по имени в главах 3, 5 и 9.
